\documentclass[11pt]{article}
\usepackage[margin=25mm]{geometry}
\usepackage{fontspec}
\setmainfont{Latin Modern Roman}
\usepackage{amsmath,amssymb,mathtools}
\usepackage{unicode-math}
\setmathfont{Latin Modern Math}
\usepackage{xurl,hyperref}
\hypersetup{colorlinks=true,urlcolor=blue}
\setcounter{secnumdepth}{0}
\setlength{\parindent}{0pt}
\setlength{\parskip}{5pt}
\setlength{\emergencystretch}{3em}
\newcommand{\sref}[2]{\hyperlink{source:#1:#2}{[S#2]}}
\newcommand{\cataloguescope}[1]{\paragraph{Recorded target.}#1}
\newcommand{\R}{\mathbb R}
\newcommand{\C}{\mathbb C}
\newcommand{\N}{\mathbb N}
\newcommand{\Z}{\mathbb Z}
\newcommand{\Q}{\mathbb Q}
\newcommand{\F}{\mathbb F}
\newcommand{\eps}{\varepsilon}
\begin{document}
% UnsolvedMath ID 30003789; source catalogue code OWR-16161-002
\setcounter{section}{30003788}
\section{Topological Defects in Three-Dimensional Polycrystals}\label{30003789}
{\small UnsolvedMath ID 30003789 · Mathematical Physics}
\subsection{Definitions and mathematical statement}

\paragraph{Shared notation.}$\mathbb N=\{1,2,\ldots\}$; the standard number fields have their usual meanings. All additional conventions are stated locally.


\paragraph{Definitions and conventions.}
A polycrystal orientation field describes local lattice orientation, typically by a map into $SO(3)$ modulo the crystal\textquotesingle s point-group symmetry. Unit quaternions give a double cover of $SO(3)$, so opposite quaternions represent the same physical rotation. Topological defects are configurations whose orientation map around a defect cannot be continuously filled within the allowed order-parameter space; line and point defects involve different homotopy tests. Symmetry quotienting matters to their classification.

\paragraph{Target retained from the canonical source catalogue.}
\paragraph{Catalogue wording retained for traceability.}
For three-dimensional orientation-field models of polycrystalline solidification, are there types of topological defects beyond those already identified in the two-dimensional and quaternionic descriptions?

\paragraph{Source scope and qualification.}
The already identified defect list and specific orientation-space quotient are not reproduced in the extraction. They must be taken from the source model; defects of an unrelated order-parameter space do not answer the question.
\subsection{Short English statement}
Are additional topological defect types present in three-dimensional orientation-field descriptions of polycrystals?
\subsection{Research attempt}
\paragraph{Model, date and investigated scope.}
OpenAI GPT-6 Astra (configured model \texttt{gpt-6-astra}), ultra reasoning; generated 11 October 2026. Verified the standard noncontractible rotation loop through quaternion lifting; no new source-model defect is identified.

\paragraph{Sources and literature check.}
Read the supplied definition, catalogue wording, source qualifications and reference entries for 30003789. The historical source/context and inherited status are retained; no fresh literature or novelty certification is claimed. No external theorem is used without proof in this attempt; explicit model assumptions are distinguished from facts about the source model.

\paragraph{A quaternion loop check.} Before looking for new defects, consider the orientation space \(SO(3)\) without a further crystal quotient. The path of unit quaternions
\[
 q(\theta)=\cos(\theta/2)+\mathbf k\sin(\theta/2),\qquad0\le\theta\le2\pi,
\]
projects to rotation by angle \(\theta\) about one fixed axis. Its projected path is a closed loop, whereas its lift runs from \(1\) to \(-1\). This loop cannot contract while its basepoint stays fixed. Indeed the two-to-one quaternion projection is locally invertible; a proposed homotopy can be lifted successively through these local inverse neighborhoods. Its terminal lift lies in the discrete fiber \(\{1,-1\}\) and varies continuously with the homotopy parameter, hence stays \(-1\). A constant loop instead has lift ending at \(1\), a contradiction.
\paragraph{Consequences and limits.} A spatial orientation field around a line can therefore carry a sign-valued obstruction invisible to a single globally chosen quaternion chart. This is the familiar rotation-loop mechanism, not a newly discovered three-dimensional defect. The source explicitly leaves the point-group quotient and already identified list to the report. Passing to \(SO(3)/G\) for a finite crystal group \(G\) changes the fiber and possible loop endpoints; one cannot simply reuse the two-point argument as a complete classification. Nor does a line-loop obstruction classify point defects, which require filling spheres. The first gap is comparison of the actual quotient's line and point classes with the report's known defects and with admissible energetic cores. No additional defect type is established here.

\paragraph{Verification and adversarial audit.}
Checked opposite-quaternion identification, continuity of the lifted endpoint, and the distinction between line loops and point-defect sphere tests. This is a GPT-6 Astra Ultra self-audit, not independent human verification. Prompt A drafting followed by prompt B scope/dependency/boundary self-audit. Checked opposite-quaternion identification, continuity of the lifted endpoint, and the distinction between line loops and point-defect sphere tests. The remaining gap is: Compute the actual crystal-quotient defect classes and compare them with the report’s identified list and admissible orientation cores. Unresolved status and zero original-target completion are retained.

Coordinator checked the lifted quaternion endpoint obstruction and the explicit limitations before imposing the actual crystal quotient. This additional check was model-assisted, not independent human verification.

\paragraph{Tools, process and contributions.}
Python 3 and local file-reading/writing tools for manuscript serialization and structural checks; mathematical calculations were symbolic reasoning, with no numerical experiment or theorem prover. Contribution prompt A was followed by prompt B. GPT-6 Astra with ultra reasoning drafted the research and performed its stated self-audit; any additional coordinator contribution is identified in the verification and audit. A GPT-6-based Codex coordinator assembled the manuscript and checked repository integration; its exact variant and reasoning configuration were not exposed. No human mathematical contribution, independent human proof check or formal proof certification is supplied. Complete derivations and the consulted source references are included here; private reasoning transcripts are not published.

\paragraph{Final status and first remaining gap.}
\textbf{UNRESOLVED}. Compute the actual crystal-quotient defect classes and compare them with the report’s identified list and admissible orientation cores.

\section{Completion Estimate}
\noindent\textbf{COMPLETION: 0\%}
The proved calculation is an elementary auxiliary result or a test of a restricted strategy; no new portion of the full source target is established.

\subsection{Sources}
\begin{itemize}
\item[\hypertarget{source:30003789:1}{S1}] ULAM.AI and the UnsolvedMath Contributors, UnsolvedMath, record 30003789 (OWR-16161-002). Catalogue curation under CC BY 4.0; definitions and English exposition added locally.
\url{https://huggingface.co/datasets/ulamai/UnsolvedMath}.
\item[\hypertarget{source:30003789:2}{S2}] Nonlinear Data: Theory and Algorithms (2018), Oberwolfach Reports, DOI 10.4171/owr/2018/20, pp. 1161--1234. Primary report accessed 2026-10-11; relevant PDF pages 18.
\url{https://doi.org/10.4171/owr/2018/20}.
\item[\hypertarget{source:30003789:3}{S3}] workshop report (catalogue-listed context; not independently consulted in this preparation).
\url{https://ems.press/content/serial-article-files/46744}.
\item[\hypertarget{source:30003789:4}{S4}] two-dimensional orientation defects (catalogue-listed context; not independently consulted in this preparation).
\url{https://doi.org/10.1103/PhysRevE.96.052802}.
\end{itemize}
\label{end:30003789}
\end{document}
